\documentclass[11pt]{article}
\usepackage[margin=1in]{geometry}
\usepackage[T1]{fontenc}
\usepackage[utf8]{inputenc}
\usepackage{lmodern}
\usepackage{amsmath}
\usepackage{enumitem}
\usepackage{amsfonts}
\usepackage{amssymb}

\setlist[itemize]{noitemsep, topsep=2pt}
\setlist[enumerate]{noitemsep, topsep=2pt}

\title{TickRange Accrual Notes Pricing}
\author{}
\date{\today}

\begin{document}
\maketitle
\tableofcontents


\section{Pricing Models}

\subsection{Model ...}

\subsubsection{Assumptions}
\subsubsection{Constrains}
\subsubsection{Variables}
- parameters
- endogeneous
- exogeneous


Having a access to a risk-free rate (say a lending protocol oracle interest rate)

Note that given \((\Delta_i)\), \(T\), \([\psi_b, \psi_a]\) And \(\Theta\):

\begin{align*}
  \Pi^{\texttt{TAN}} (T) &= \Theta \cdot \frac{\sum^N \mathbb{1}_{\psi_b \leq P(\Delta_i) \leq \psi_a}}{N} \\
  \\
  &= \frac{\Theta}{N}\cdot \sum \mathbb{1}_{P(\Delta_i) > \psi_b} - \sum \mathbb{1}_{P(\Delta_i) > \psi_a}
\end{align*}

The above decomposition is a series of digital call options (cash-or noting calls)


Then considering a period tick accrual policy (times when tick accrual can accumulate coupons)
\(0 \leq T_1 \leq \cdots \leq T_N = T\). Then using the martinagle argument, we have:

\begin{align*}
  P_{\Pi^{\text{TAN}}} (0, T) = \e^{-rT} \frac{\Theta}{N}\cdot \sum^N \e^{r\cdot T_i} \big ( P_{Pi^D} (0, T_i, \psi_b) - P_{\Pi^{D}} (0, T_i , \psi_a)\big)

\end{align*}

\subsection{Calibration}
\subsection{Empirical Analysis}
\section{References}

- cite the Range Accrual Strucutred NOTes by the Bulgarian thesis
\end{document}
